\chapter{Clase 17: Solución analítica del problema relativo y plano de Laplace}

Esta clase continúa el cierre analítico del problema relativo de dos cuerpos. A partir de las cuadraturas obtenidas en la clase anterior, se reduce el problema al plano orbital y se deriva la ecuación polar de la trayectoria. El objetivo central es mostrar que la solución relativa corresponde siempre a una sección cónica y que su geometría queda completamente determinada por las constantes del movimiento.

\section{Reducción al plano de la trayectoria}
Partimos del conteo de cuadraturas obtenido en el espacio tridimensional:
\begin{itemize}
    \item momento angular relativo específico \(\vec{h}\): 3 componentes,
    \item energía relativa específica \(\mathcal{E}\): 1 escalar,
    \item vector de Laplace o de excentricidad \(\vec{e}\): 3 componentes.
\end{itemize}

En apariencia esto produce 7 constantes, pero una de ellas queda ligada por la relación
\[
\vec{h} \cdot \vec{e} = 0,
\]
por lo que el sistema tiene finalmente \textbf{6 cuadraturas independientes}.

Al pasar al \textbf{plano de la trayectoria} —también llamado plano orbital o plano perpendicular a \(\vec{h}\)— el problema se reduce a dos dimensiones espaciales. En ese plano, el número de variables dinámicas independientes es \(4\): \((x, y, \dot{x}, \dot{y})\). En consecuencia, basta con describir el movimiento mediante cuatro constantes escalares equivalentes:
\begin{enumerate}
    \item \(h = |\vec{h}|\), módulo del momento angular específico,
    \item \(\mathcal{E}\), energía relativa específica,
    \item la dirección de \(\vec{e}\) en el plano orbital,
    \item el módulo \(e = |\vec{e}|\).
\end{enumerate}

Este marco permite enfocarse exclusivamente en la geometría y dinámica plana de la órbita.

\begin{importante}
La reducción al plano orbital no es una aproximación. Es una consecuencia exacta de la conservación del momento angular: toda la dinámica relativa ocurre en un único plano fijo.
\end{importante}

\section{Análisis del vector de Laplace}
El vector de excentricidad se define como:
\[
\vec{e} = \frac{\vec{v} \times \vec{h}}{\mu} - \frac{\vec{r}}{r}
\]

\subsection{Dirección y significado geométrico}
Sus propiedades más importantes son:
\begin{itemize}
    \item \(\vec{e}\) es constante en magnitud y dirección,
    \item está contenido en el plano orbital, pues es combinación lineal de \(\vec{r}\) y \(\vec{v}\),
    \item apunta hacia el \textbf{periastro}, es decir, el punto de mínima distancia radial.
\end{itemize}

El ángulo \(\theta\) medido desde la dirección de \(\vec{e}\) hasta \(\vec{r}\) se denomina \textbf{verdadera anomalía}. Esta elección angular es especialmente natural, porque toma como referencia el eje geométricamente privilegiado de la órbita.

\subsection{Cálculo de la magnitud \texorpdfstring{$e^2$}{e²}}
Calculamos el producto escalar \(\vec{e} \cdot \vec{e}\) para obtener una expresión analítica en función de las integrales del movimiento:
\[
 e^2 = \left( \frac{\vec{v} \times \vec{h}}{\mu} - \frac{\vec{r}}{r} \right) \cdot \left( \frac{\vec{v} \times \vec{h}}{\mu} - \frac{\vec{r}}{r} \right)
\]

Desarrollando:
\[
 e^2 = \frac{|\vec{v} \times \vec{h}|^2}{\mu^2} - \frac{2}{\mu r}(\vec{v} \times \vec{h}) \cdot \vec{r} + 1
\]

\textbf{Paso A. Primer término: identidad de Lagrange}
\[
 |\vec{v} \times \vec{h}|^2 = v^2 h^2 - (\vec{v} \cdot \vec{h})^2
\]
Como \(\vec{h} = \vec{r} \times \vec{v}\), entonces \(\vec{v} \cdot \vec{h} = 0\). Por tanto:
\[
 |\vec{v} \times \vec{h}|^2 = v^2 h^2
\]

\textbf{Paso B. Segundo término: propiedad cíclica del triple producto escalar}
\[
 (\vec{v} \times \vec{h}) \cdot \vec{r} = \vec{h} \cdot (\vec{r} \times \vec{v}) = \vec{h} \cdot \vec{h} = h^2
\]

\textbf{Paso C. Ensamblaje y sustitución de la energía}
\[
 e^2 = \frac{v^2 h^2}{\mu^2} - \frac{2 h^2}{\mu r} + 1
\]
\[
 e^2 = 1 + \frac{2h^2}{\mu^2}\left(\frac{v^2}{2} - \frac{\mu}{r}\right)
\]

Reconociendo que
\[
\mathcal{E} = \frac{v^2}{2} - \frac{\mu}{r},
\]
obtenemos la relación fundamental:
\[
\boxed{e^2 = 1 + \frac{2\mathcal{E}h^2}{\mu^2}}
\]

\begin{importante}
Esta ecuación conecta directamente la energía del sistema con la forma geométrica de la órbita. La excentricidad no es un parámetro puramente geométrico aislado: está determinada por las constantes dinámicas del movimiento.
\end{importante}

\subsection{Clasificación de órbitas según la energía}
La relación anterior permite clasificar las cónicas orbitales:
\begin{itemize}
    \item si \(\mathcal{E} < 0\), entonces \(0 \le e < 1\): la órbita es una \textbf{elipse},
    \item si \(\mathcal{E} = 0\), entonces \(e = 1\): la órbita es una \textbf{parábola},
    \item si \(\mathcal{E} > 0\), entonces \(e > 1\): la órbita es una \textbf{hipérbola}.
\end{itemize}

Así, la energía específica determina si la trayectoria es ligada, marginal o de escape.

\section{Deducción de la ecuación polar de la trayectoria}
Proyectamos ahora el vector de Laplace sobre el vector posición \(\vec{r}\) para obtener la relación funcional entre \(r\) y \(\theta\).

\subsection{Proyección geométrica}
Por definición del producto punto:
\[
\vec{e} \cdot \vec{r} = e r \cos\theta
\]
donde \(\theta\) es la verdadera anomalía, es decir, el ángulo entre \(\vec{e}\) y \(\vec{r}\).

\subsection{Proyección algebraica}
Usando la definición del vector \(\vec{e}\):
\[
\vec{e} \cdot \vec{r} = \left(\frac{\vec{v} \times \vec{h}}{\mu} - \frac{\vec{r}}{r}\right) \cdot \vec{r}
\]
\[
\vec{e} \cdot \vec{r} = \frac{(\vec{v} \times \vec{h}) \cdot \vec{r}}{\mu} - \frac{\vec{r} \cdot \vec{r}}{r}
\]

Usando que \((\vec{v} \times \vec{h}) \cdot \vec{r} = h^2\) y que \(\vec{r} \cdot \vec{r} = r^2\), resulta:
\[
\vec{e} \cdot \vec{r} = \frac{h^2}{\mu} - r
\]

\subsection{Igualación y despeje de \texorpdfstring{$r(\theta)$}{r(theta)}}
Igualando la expresión geométrica con la algebraica:
\[
 e r \cos\theta = \frac{h^2}{\mu} - r
\]
\[
 r(1 + e\cos\theta) = \frac{h^2}{\mu}
\]

Definiendo el \textbf{parámetro} o \textbf{semilatus rectum}
\[
 p \equiv \frac{h^2}{\mu},
\]
se llega a la ecuación polar final de la trayectoria:
\[
\boxed{r(\theta) = \frac{p}{1 + e\cos\theta}}
\]

Esta es la ecuación general de una sección cónica con un foco en el origen del sistema relativo.

\begin{importante}
La ecuación \(r(\theta) = \frac{p}{1 + e\cos\theta}\) muestra que el problema relativo de dos cuerpos tiene como soluciones exactas cónicas: circunferencias, elipses, parábolas e hipérbolas aparecen como casos particulares según el valor de \(e\).
\end{importante}

\section{Resumen de la integración completa}
Los resultados principales de esta clase pueden sintetizarse así:
\begin{enumerate}
    \item El problema relativo de dos cuerpos posee suficientes cuadraturas independientes para ser completamente integrable.
    \item En el plano orbital, la dinámica queda caracterizada por \(h\), \(\mathcal{E}\), la dirección de \(\vec{e}\) y su magnitud \(e\).
    \item El vector \(\vec{e}\) fija la orientación de la órbita y su módulo determina la excentricidad geométrica mediante
\[
 e^2 = 1 + \frac{2\mathcal{E}h^2}{\mu^2}.
\]
    \item La proyección de \(\vec{e}\) sobre \(\vec{r}\) conduce directamente a
\[
 r(\theta) = \frac{p}{1 + e\cos\theta},
\]
lo que demuestra analíticamente que la trayectoria relativa es siempre una cónica.
\end{enumerate}

Este desarrollo cierra la solución analítica del problema de dos cuerpos y deja listo el marco conceptual para estudiar con más detalle la geometría orbital, los elementos de la cónica y su representación en distintos sistemas de referencia.